\begin{figure}[p]
\centering
\begin{tikzpicture}[gclplate,scale=.86]
\node[anchor=west,font=\sffamily\bfseries\Large,gcllabel] at (-6,3.9) {A NORMALIZATION TRACE};
\draw[GCLWarm,line width=1pt] (-6,3.55)--(6,3.55);
\node[gcllabel] (s0) at (0,2.35) {$((\lambda f{:}P\to P.\lambda y{:}P.f\,y)(\lambda x{:}P.x))$};
\node[gcllabel] (s1) at (0,1.15) {$\lambda y{:}P.(\lambda x{:}P.x)\,y$};
\node[gcllabel] (s2) at (0,-.05) {$\lambda y{:}P.y$};
\draw[gclbluearrow] (s0)--(s1);
\draw[gclbluearrow] (s1)--(s2);
\node[gcllabel,font=\bfseries] at (4.8,2.35) {$P\to P$};
\node[gcllabel,font=\bfseries] at (4.8,1.15) {$P\to P$};
\node[gcllabel,font=\bfseries] at (4.8,-.05) {$P\to P$};
\draw[GCLGreen,line width=1pt] (4.15,2.75)--(4.15,-.45);
\node[gcllabel,rotate=90] at (5.55,1.15) {type invariant};
\node[gcllabel,align=center] at (0,-1.0) {trace generated by \texttt{lab02\_normalizer.py}};
\end{tikzpicture}
\caption{\textbf{Plate 4. A normalization trace.} The teaching normalizer contracts two beta redexes and rechecks the invariant type after each step. This computed trace is evidence for the implementation on the retained fixture; one trace is not a proof of strong normalization.}
\label{plate:normalization-trace}
\end{figure}
